\chapter{Naming the material, export, and using Indradhanu V2 in a chain}\label{ch:material}

\begin{onething}
A score map ranks pixels; a person turns the ranking into a decision by choosing a percentile, the flagged
pixels are then compared with a library of laboratory spectra to propose what they are made of, and the
result leaves the system as a bundle whose coordinates are recovered, at the last moment, from the raw file.
Running the newer model, Indradhanu~V2, in that chain needs a band filter that keeps native bands, the
default since 2026-10-31. A redesign of the naming step, agreed and built on 2026-10-31 (\cref{ch:namer}),
asks first whether nature could explain a pixel and only then which man-made class it is.
\end{onething}

\Cref{ch:inference} reconstructed a whole scene and \cref{ch:residualmaha} turned the residual into a score
per pixel. The product's chain does not stop there (\cref{fig:inference:chain}). Three steps remain, and this
chapter takes them in order: the threshold, which only a person may set; the naming of materials, whose
mathematics fills a companion volume and is only sketched here; and the export, where geography returns.
The chapter ends with the analyst's view of Indradhanu~V2: what to change in the chain to use it, what goes
wrong if you do not, and what extra files it leaves behind.

%======================================================================
\section{Why a threshold is a person's job}\label{sec:material:why}

You have a score map. Which pixels are anomalies? The obvious fix, a fixed cut on the score, cannot work,
because typical scores differ by an order of magnitude between scenes (AD12; the numbers are
illustrative \tagA):
\begin{center}\small
\begin{tabular}{@{}lccl@{}}
\toprule
scene & typical score & fixed cut & result\\
\midrule
calm lake & 0.01 & 0.05 & almost nothing flagged\\
burned area & 0.10 & 0.05 & most of the scene flagged\\
\bottomrule
\end{tabular}
\end{center}
The same number is far too strict for one scene and far too loose for the other. Residual Mahalanobis makes
this worse, not better: it is a squared distance whose scale depends on the scene's own covariance and on how
many directions were kept.

So the convention is a \emph{percentile}\index{percentile threshold}: ``flag the top $1\%$ of this scene''.
That adapts to each scene's own distribution, as a person looking at the map would. And \emph{which}
percentile is not a property of the data at all: it depends on how many candidates the team can review, how
costly a miss is, and what the scene is. That is why the action that prepares the decision,
\code{anomaly_detection_prep}, ends in a status no other action uses, \code{needs_threshold}, and waits for a
person (\file{backend/allotrope/action_types/anomaly_detection_prep.py}, line 53).

%======================================================================
\section{From score maps to an anomaly mask}\label{sec:material:mask}

\subsection{Rescale, weight, combine}

\code{anomaly_detection_prep} reads every model's \file{anomaly_score.tif} from one upstream
\code{anomaly_scoring} output and builds one \emph{composite score}\index{composite score}
(\file{backend/allotrope/action_types/_anomaly_detection_prep_run.py}, lines 88--191). Three steps:
\begin{enumerate}
\item \textbf{Min--max rescale each map} over its finite pixels,
  \begin{equation}
  \tilde{s}(p)=\frac{s(p)-s_{\min}}{s_{\max}-s_{\min}}\in[0,1],
  \end{equation}
  so that maps in different units (squared distances, angles, reflectance) become comparable. A flat map
  becomes all zeros.
\item \textbf{Normalise the weights}: the user's \code{algorithm_weights} (equal by default) are divided by
  their sum, $\omega_m=w_m/\sum_k w_k$, so the composite stays in $[0,1]$.
\item \textbf{Combine}: $c(p)=\sum_m\omega_m\tilde{s}_m(p)$. A pixel is finite in the composite only if it is
  finite in every contributing map; otherwise it is NaN (not a number).
\end{enumerate}
Note that this is a different normalisation from \texttt{combined}'s, which divides by the maximum only
(\cref{we:inference:combined}).

\begin{worked}[label={we:material:composite}]{A two-model composite at one pixel}
Indradhanu's map runs from 0 to 36 and MNF-RX's from 2 to 82 over the scene; at one pixel they read 7.2 and
30 \tagA. The user gives Indradhanu weight 2 and MNF-RX weight 1. On paper:
\begin{align*}
\tilde{s}_{\text{Indradhanu}} &= \frac{7.2-0}{36-0} = \frac{7.2}{36} = 0.20,\\
\tilde{s}_{\text{MNF-RX}} &= \frac{30-2}{82-2} = \frac{28}{80} = 0.35,\\
\omega_{\text{Indradhanu}} &= \frac{2}{2+1} = \frac23, \qquad \omega_{\text{MNF-RX}} = \frac{1}{2+1} = \frac13,\\
c &= \tfrac23\times0.20+\tfrac13\times0.35\\
&= 0.1333+0.1167\\
&= 0.25.
\end{align*}
\end{worked}

\subsection{The percentile, as the code computes it}

The viewer shows a slider. Each press of \emph{Apply} calls
\code{POST /actions/\{id\}/anomaly_detection_preview}, which converts the slider value $q$ (a percentile,
0 to 100) into an absolute threshold and binarises the composite
(\file{backend/allotrope/api/_anomaly_detection_preview.py}, lines 168--208):
\begin{lstlisting}
absolute_threshold = float(np.percentile(finite_values, q))   # finite_values: sorted, (N,)
raw_mask = composite >= np.float32(absolute_threshold)         # (H, W) bool
raw_mask &= valid_mask                                         # NaN pixels are never flagged
\end{lstlisting}
An optional square dilation\index{dilation} then grows the mask by a user-chosen odd kernel, and if a ground-truth annotation
is attached the endpoint reports precision, recall and F1 (\cref{ch:measures}). Nothing is written until the
user presses \emph{Commit}, which calls \code{POST /actions/\{id\}/anomaly_detection_commit} and writes
\file{anomaly_mask.tif} and \file{metrics.json} (\file{backend/allotrope/api/action_threshold.py}, lines
283--330).

\begin{sourcenote}[: commit is an endpoint, not an action]
The prep action's docstring says finalisation ``is deferred to a separate \texttt{anomaly\_detection\_commit}
action'' (\file{anomaly_detection_prep.py}, lines 21--25). In the source, commit is an api endpoint on the
prep action itself, \code{anomaly_detection_commit} in \file{backend/allotrope/api/action_threshold.py}; there
is no such action type in \file{backend/allotrope/action_types/}.
\end{sourcenote}

\paragraph{How many pixels does a percentile flag?} \code{np.percentile} with its default (linear)
interpolation places the $q$-th percentile of $N$ sorted values at position
\begin{equation}
\text{pos}=\frac{q}{100}\,(N-1)
\end{equation}
in the list, counting from 0, interpolating between neighbours when \text{pos} is not whole. If the values are
all different, the ones at or above that threshold are those at positions $\lceil\text{pos}\rceil,\dots,N-1$,
so the number flagged is
\begin{equation}
n_{\text{flag}}=N-\lceil\text{pos}\rceil.
\label{eq:material:nflag}
\end{equation}

\begin{worked}[label={we:material:percentile}]{Flagging the top 30\% of ten pixels, and the top 0.5\% of a scene}
Ten composite values, sorted: $0.02, 0.05, 0.07, 0.10, 0.12, 0.15, 0.21, 0.33, 0.48, 0.90$ \tagA. The slider is
at $q=70$:
\begin{align*}
\text{pos} &= \frac{70}{100}\times(10-1) = 0.7\times9 = 6.3.
\end{align*}
Position 6 holds $0.21$ and position 7 holds $0.33$; interpolate $0.3$ of the way:
\begin{align*}
\text{threshold} &= 0.21+0.3\times(0.33-0.21) = 0.21+0.3\times0.12 = 0.21+0.036 = 0.246.
\end{align*}
Flagged: the values $\ge0.246$, which are $0.33, 0.48, 0.90$: three pixels, and
\cref{eq:material:nflag} agrees, $10-\lceil6.3\rceil=10-7=3$.

Now the scene of \cref{we:inference:funnel}, with $N=956{,}679$ kept pixels and $q=99.5$:
\begin{align*}
\text{pos} &= 0.995\times(956{,}679-1) = 0.995\times956{,}678 = 951{,}894.61,\\
n_{\text{flag}} &= 956{,}679-\lceil951{,}894.61\rceil = 956{,}679-951{,}895 = 4{,}784.
\end{align*}
The course's ``$956{,}679\times0.005=4{,}783$'' rounds the other way; the code flags 4,784 (checked by the
book's script on distinct values \tagC). Ties, which are common in maps with many exact zeros, can push the
count higher, because every pixel equal to the threshold is flagged.
\end{worked}

\subsection{Which pixels the percentile counts}

Whether ``$N$'' above is the kept pixels or the whole scene depends on the upstream models, because of how
each kind marks the pixels it did not score.
\begin{itemize}
\item A \emph{classical} detector writes NaN outside its mask (\file{_anomaly_scoring_run.py}, lines
  535--541, the comment on classical NaNs). NaN pixels are not finite, so they are left out of the percentile
  population.
\item A \emph{foundation} model's score is set to exactly $0$ outside the keep mask
  (\file{app/utils/anomaly_detection/scoring.py}, line 114), and the GeoTIFF is written with no nodata value.
  Zero is finite, so those pixels \emph{are} in the population, at the bottom.
\end{itemize}

\begin{sourcenote}[: a model-only chain takes the percentile over the whole frame]
The anomaly course's end-to-end walk (AD13) flags ``the top 0.5\%'' of the 956,679 kept pixels. That holds
when a classical detector (MNF-RX) is in the composite, because its NaNs remove the unkept pixels from the
population. With Indradhanu alone, every one of the $1{,}474{,}990$ pixels of the frame is finite, the
$518{,}311$ unkept ones being exact zeros, so $q=99.5$ flags
\begin{align*}
1{,}474{,}990-\lceil0.995\times1{,}474{,}989\rceil = 1{,}474{,}990-1{,}467{,}615 = 7{,}375 \text{ pixels},
\end{align*}
which is $7{,}375/956{,}679=0.771\%$ of the kept pixels, not $0.5\%$ \tagC. To flag about 4,784 the slider would
have to be at $q\approx99.676$. The cut is still applied to a ranking the analyst can see, so nothing is
wrong; but ``top 0.5\%'' means 0.5\% of the finite pixels, and which pixels are finite depends on the models
chosen. Sources: \file{_anomaly_detection_preview.py}, line 84 (\code{valid_mask = np.isfinite}), and
\file{_anomaly_detection_prep_run.py}, lines 179--191.
\end{sourcenote}

%======================================================================
\section{Naming the material}\label{sec:material:naming}

\subsection{The question, and why only hyperspectral can answer it}

\begin{sloppypar}
A shortlist of unusual pixels is not yet useful. An analyst wants ``probably a metal roof'' or ``probably
calcite''. Hyperspectral data can propose that, because a spectrum is close to a material's fingerprint
(\cref{ch:light}); a single thermal band cannot. The step that does it is the action
\code{spectral_library_match}, implemented in
\file{backend/allotrope/action_types/_spectral_library_match_run.py} on top of \file{app/spectral_match/}.
Its full mathematics (resampling a laboratory spectrum through each band's response, the spectral angle,
whitened target detection, cones of mixtures, non-negative least squares and the calibration against a
scene's ordinary pixels) is the subject of \emph{the companion volume on spectral matching}
(\file{research/math/spectral-matching/}). This section gives the intuition and the outputs, so that you can
read the chain end to end.
\end{sloppypar}

\subsection{It reads the native cube, not the band-filtered one}

The action loads the \emph{native onboarding vendable}, not \code{band_filter_apply}'s output
(\file{_spectral_library_match_run.py}, lines 155--164). Resampling onto the 165-band common grid
(\cref{ch:bandpipeline}) is what makes sensors comparable for a model, but it blunts narrow absorption
features, and those features are the evidence a material match relies on.

\subsection{The spectral angle against a library, in one example}

The library is USGS splib07 \citep{kokaly2017usgs}, curated to 519 laboratory spectra in
\file{data/splib07_slim/} \tagM{docs/training/anomaly-detection/14-material-and-export.md}. Each library spectrum
is first simulated through each band of the sensor: a band is not a point measurement but a weighted
average over a range of wavelengths, roughly Gaussian with full width at half maximum (FWHM)
$\text{FWHM}=2\sqrt{2\ln2}\,\sigma\approx2.3548\,\sigma$. For a 10~nm band,
\begin{align*}
\sigma &= \frac{10}{2.3548} = 4.247\text{ nm}.
\end{align*}
The first matcher, the \emph{spectral angle mapper}\index{spectral angle mapper} \citep{kruse1993spectral},
then ranks library entries by the angle between the pixel and each entry (\cref{ch:losses}). The angle ignores
brightness, which is what you want for sun versus shadow, but it has a basic problem at 30-metre pixels: a
roof, a truck or a tarpaulin rarely fills a pixel.

\begin{worked}[label={we:material:mixed}]{SAM names the grass around a roof}
The anomaly course's three-band example (AD15): grass $\vect{g}=(0.05,0.40,0.25)$, roof
$\vect{r}=(0.30,0.32,0.35)$, and a pixel that is 80\% grass and 20\% roof:
\begin{align*}
x_1 &= 0.8\times0.05+0.2\times0.30 = 0.040+0.060 = 0.100,\\
x_2 &= 0.8\times0.40+0.2\times0.32 = 0.320+0.064 = 0.384,\\
x_3 &= 0.8\times0.25+0.2\times0.35 = 0.200+0.070 = 0.270.
\end{align*}
Angle to grass:
\begin{align*}
\vect{x}\cdot\vect{g} &= 0.100\times0.05+0.384\times0.40+0.270\times0.25 = 0.0050+0.1536+0.0675 = 0.2261,\\
\norm{\vect{x}} &= \sqrt{0.0100+0.147456+0.0729} = \sqrt{0.230356} = 0.47995,\\
\norm{\vect{g}} &= \sqrt{0.0025+0.1600+0.0625} = \sqrt{0.2250} = 0.47434,\\
\cos\theta_g &= \frac{0.2261}{0.47995\times0.47434} = \frac{0.2261}{0.22766} = 0.99315,
\quad\theta_g = 6.7^\circ.
\end{align*}
Angle to roof:
\begin{align*}
\vect{x}\cdot\vect{r} &= 0.0300+0.12288+0.0945 = 0.24738,\\
\norm{\vect{r}} &= \sqrt{0.09+0.1024+0.1225} = \sqrt{0.3149} = 0.56116,\\
\cos\theta_r &= \frac{0.24738}{0.47995\times0.56116} = \frac{0.24738}{0.26933} = 0.91851,
\quad\theta_r = 23.3^\circ.
\end{align*}
SAM confidently names the surroundings. Subtracting the background first fixes it: with a ring of pure grass
around the pixel, $\vect{x}-\vect{g}=(0.050,-0.016,0.020)$, which is exactly $0.2\times(\vect{r}-\vect{g})$.
The difference points along ``roof minus grass''; the 80\% grass has vanished.
\end{worked}

That subtraction is the first of five ideas in the \emph{family matcher}\index{family matcher}
(\file{app/spectral_match/family_match.py}) that has run after SAM since 2026-10-16 (AD15), and whose answer the viewer and the export now lead with. In one line each:
\begin{enumerate}
\item subtract the mean of a ring around the pixel ($13\times13$ minus the central $5\times5$, 144 pixels;
  the same ring idea as \cref{ch:residualmaha}, smaller);
\item measure how well the difference points along ``library spectrum minus ring'', after whitening, with the
  adaptive coherence estimator;
\item score a whole \emph{family} of spectra (roofing, metal, plastic, \dots) as any non-negative mix of its
  members, found by non-negative least squares, for man-made families only;
\item calibrate each family score as a percentile against up to 10,000 ordinary pixels of the same scene;
\item give a man-made name only if that percentile reaches $p^*=1-\alpha/F$, where $\alpha$ is the share of
  ordinary pixels you accept being given a man-made name and $F$ the number of man-made families.
\end{enumerate}
With the defaults, $\alpha=0.05$ (\code{DEFAULT_ALPHA}, \file{family_match.py}, line 36) and $F=10$ man-made
families in the shipped library \tagM{docs/training/anomaly-detection/15-naming-the-family.md}:
\begin{align*}
\frac{\alpha}{F} &= \frac{0.05}{10} = 0.005,\\
p^* &= 1-0.005 = 0.995.
\end{align*}
By the union bound\index{union bound}, an ordinary pixel then gets \emph{some} man-made name with probability at most
$10\times0.005=0.05=\alpha$. Why this works, and what it does not promise, is worked in full in the companion
volume.

\subsection{What the action writes}

\begin{center}\small
\begin{tabular}{@{}lp{0.62\linewidth}@{}}
\toprule
file & contents\\
\midrule
\texttt{matches.parquet} & SAM: top-$K$ library entries and angles per flagged pixel\\
\texttt{match\_map.tif} & SAM: top-1 library index per pixel; $-2$ outside the mask, $-1$ no match\\
\texttt{family\_matches.parquet} & family, group, calibrated confidence, named or not, best single entry\\
\texttt{family\_map.tif} & int16 family code per pixel (below)\\
\texttt{family\_legend.json} & code $\to$ family, group, colour; $\alpha$ and $p^*$\\
\texttt{summary.json} & counts per family, timings, \texttt{family.status}\\
\bottomrule
\end{tabular}
\end{center}
The family codes (\file{_spectral_library_match_run.py}, lines 53--56 and 396--399):
\begin{center}\small
\begin{tabular}{@{}cl@{}}
\toprule
code & meaning\\
\midrule
$-2$ & not an anomaly pixel\\
$-1$ & not scored (nodata on a band the matcher uses)\\
$0$ & best family is man-made but below $p^*$: ``man-made, type unconfirmed''\\
$1,\dots,17$ & that family, in sorted name order, plus one (natural families are shown at any confidence)\\
\bottomrule
\end{tabular}
\end{center}
Family matching never fails the action: an exception is recorded as \code{family.status = "failed"} in the
summary, and the SAM outputs still make a complete result. It runs only in \code{mode="anomaly_pixels"},
because its per-pixel least squares cost about 2.5~ms per pixel (45~s for about 18,000 pixels, AD15): a
million pixels would take about $1{,}000{,}000\times0.0025\text{ s}=2{,}500\text{ s}\approx42$ minutes.

\begin{keypoint}
The output is a ranked hypothesis about a pixel the analyst already kept, not a filter that decides which
pixels are real: false alarms outnumber targets about 50 to 1 on the PRISMA scenes, so most man-made names in
a scene still land on false alarms (AD15).
\end{keypoint}

%======================================================================
\section{A redesign: natural first, then the class}\label{sec:material:twolevel}

\begin{sourcenote}[: the design; the model is in \cref{ch:namer}]
This section is the \emph{why} of a redesign of the naming step agreed between the user and Claude on 2026-10-31.
It was written up the same day as \file{docs/proposals/spectral-namer.md} and \emph{built} the same day in
\file{app/spectral_match/namer/}; its first run, N001, started training on an H100 at 11:30 UTC and has no results
yet. \Cref{ch:namer} teaches the model as built, which departs from this section's design in places, most of all in
how level~1 is trained (the note after \cref{we:material:openset}). Nothing here runs in the product: the namer is
wired into no action until it passes. The experiment README's heading ``Design so far'' predates the proposal.
The one classical measurement, the separability check of \cref{sec:material:separability}, is
\file{research/experiments/spectral-namer/separability.py}. The companion volume's Part~VII derives the
mathematics slowly.
\end{sourcenote}

\subsection{Where the ten man-made families came from}

\file{scripts/curate_family_library.py} builds \file{data/family_library/} from the USGS and ECOSTRESS libraries
and gives every spectrum a group and a family by \emph{regular expressions}\index{regular expression} (text
patterns) on its file name, read top to bottom, first match wins. Four of the USGS chapter-A rules, in order:
\begin{lstlisting}
(r"asphalt_shingle|asphalt_tar|tar_paper|roofing_felt|roofng|roofing|clear_roof|cedar_shake", "manmade", "roofing"),
(r"asphalt|concrete_gds375|brick.*pav", "manmade", "pavement"),
(r"brick|cinder|concrete", "manmade", "building"),
(r"plastic|polystyrene|tyvek|visqueen|fiberglass", "manmade", "plastic"),
\end{lstlisting}
So an asphalt shingle is roofing and an asphalt road is pavement; \texttt{concrete\_gds375} is pavement and every
other concrete is building. Three consequences:
\begin{enumerate}
\item \textbf{Families by use, not material.} Asphalt sits in two families and concrete in two, decided by a word in
  a file name. From orbit a sensor sees what a surface is made of, not what it is for.
\item \textbf{Family sizes reflect the libraries.} From \file{index.json}: building 11, glass 1, metal 13, oil 40,
  paint 8, pavement 15, plastic 39, roofing 27, textile 14, wood\_paper 8; the 40s are \code{FAMILY_CAP = 40}
  (\file{app/spectral_match/family_library.py}). Man-made $11+1+13+40+8+15+39+27+14+8=176$, natural $235$, total
  $411$.
\item \textbf{Some families cannot be told apart.} The smallest principal angle between the rank-3 subspaces of
  pavement and building is $0.51^\circ$; metal and plastic $2.04^\circ$; paint and plastic $5.92^\circ$ (companion
  volume, \texttt{tab:angles}, computed there). At $0.51^\circ$ the share of a unit target's energy outside the
  other family is $\sin^2(0.51^\circ)=0.0089^2=0.00008$: one material with two names.
\end{enumerate}

\subsection{What can be named from orbit}

A car of $4.5\times1.8=8.1$\,m$^2$ fills $8.1/900=0.009$ of a 30\,m pixel and $8.1/25=0.32$ of a 5\,m one. At
PRISMA's and EnMAP's 30\,m only large objects can be named; AVIRIS-NG, flown low, reaches vehicle scale (about
5\,m in the discussion; it depends on flight height \tagA). The discussion replaced the use-based families with a
defence-oriented class list: \emph{camouflage} and \emph{military coatings} (no library spectra at all; they need
field spectra), \emph{painted metal}, \emph{bare metal}, \emph{vessels on water}, \emph{synthetic fabric and
tarpaulins}, \emph{fuel and oil}, \emph{concrete and asphalt} (one class, ending the split above), \emph{disturbed
earth}, \emph{cut vegetation}, \emph{glass and photovoltaic panels}, and \emph{none / natural}. The two classes an
analyst wants most are the two no library holds.

\subsection{The two-level model}
\index{two-level namer}

Naming becomes two questions in turn. Level~1: could nature explain this pixel? Level~2, only if not: which
man-made class? There are three possible answers: \emph{natural}, \emph{man-made, class unknown}, and \emph{man-made
plus a class}. Write $M$ for ``man-made'' and $k$ for a class. Because a class can only be man-made, the event
``class $k$'' is the event ``man-made and class $k$'', and the definition of conditional probability gives
\begin{equation}
P(k)=P(M)\,P(k\mid M).
\label{eq:material:twolevel}
\end{equation}

\begin{worked}[label={we:material:twolevel}]{The product of the two levels}
Level~1 gives $P(M)=0.9$; level~2 gives painted metal 0.6, bare metal 0.25, tarp 0.15 \tagA.
\begin{align*}
P(\text{painted metal}) &= 0.9\times0.6 = 0.54\\
P(\text{bare metal}) &= 0.9\times0.25 = 0.225\\
P(\text{tarp}) &= 0.9\times0.15 = 0.135\\
P(\text{natural}) &= 1-0.9 = 0.10\\
\text{sum} &= 0.54+0.225+0.135+0.10 = 1.00 .
\end{align*}
With $P(M)=0.3$ instead, painted metal is $0.3\times0.6=0.18$. The same level-2 answer is worth 0.54 behind a
confident level~1 and 0.18 behind a doubtful one: the product stops a confident class name from hiding a doubtful
``is it man-made at all?''.
\end{worked}

\paragraph{Level 1 is a model of nature.} The obvious level~1 is a classifier trained on library spectra of both
kinds. Its man-made side would only know library materials, and camouflage is not in the library. The natural side
is different: every scene is mostly natural pixels, and the project holds billions of them (SpectralEarth and the
native shards). So level~1 is trained to say whether nature explains a pixel, and anything it cannot explain is
called man-made, library or not. This is \emph{open-set}\index{open-set recognition} recognition: the model may meet
something it has never seen and say so.

\begin{worked}[label={we:material:openset}]{Near the natural average, off the natural line}
In two dimensions, let nature lie on the line $x+y=1$, sampled at $(0.1,0.9)$, $(0.3,0.7)$, $(0.5,0.5)$, $(0.7,0.3)$,
$(0.9,0.1)$; the library has one man-made spectrum at $(0.9,0.9)$; a camouflage net sits at $(0.4,0.4)$ \tagA.
\begin{align*}
\text{natural average} &= (0.5,\ 0.5)\\
\text{camouflage to it} &= \sqrt{0.1^2+0.1^2}=\sqrt{0.02}=0.141\\
\text{camouflage to the library} &= \sqrt{0.5^2+0.5^2}=\sqrt{0.5}=0.707\\
&\Rightarrow\ \text{``natural'' by nearest average;}\\
\text{residual } |x+y-1| \text{ of each natural sample} &= 0\\
\text{residual of the camouflage} &= |0.4+0.4-1|=0.2\\
&\Rightarrow\ \text{not natural.}
\end{align*}
The camouflage is closer to the natural average than the two end samples are (each $\sqrt{0.4^2+0.4^2}=0.566$ from
it), so no rule built on distance to an average can catch it. It is off the line nature lives on, and a model of that
line notices.
\end{worked}

\begin{sourcenote}[: as built, level 1 is a two-class head]
The built namer does not train level~1 as a model of nature. Its level-1 logit is trained with binary cross-entropy
on planted samples, library \emph{man-made} spectra mixed into real pixels against natural library spectra and plain
pixels (\file{app/spectral_match/namer/train.py}, \code{loss}; \file{generator.py}, the three sample kinds). It
does see man-made examples. Whether it still flags a material it never saw is what run N002 measures, by holding the
whole paint/coatings class out (\cref{ch:namer}).
\end{sourcenote}

\subsection{How ACE matches, and why the fill drops out}

Level~2 needs a match score. The separability check uses ACE in the form of \cref{sec:material:naming}:
$\vect{d}=\vect{x}-\vect{m}$ (pixel minus ring mean), $\vect{s}_j=\vect{t}_j-\vect{m}$ (library spectrum $j$ minus ring
mean), both whitened, and
\begin{equation}
\text{ACE}_j=\frac{(\vect{s}_j\cdot\vect{d})^2}{(\vect{s}_j\cdot\vect{s}_j)(\vect{d}\cdot\vect{d})},
\end{equation}
the squared cosine between the two whitened directions.

\begin{worked}[label={we:material:ace}]{ACE on three bands, at two fills}
Ring mean grass $(0.05,0.40,0.25)$, roof $(0.30,0.32,0.35)$, soil $(0.20,0.25,0.30)$, band standard deviations
$(0.02,0.04,0.01)$ \tagA. Whitened directions (difference divided band by band):
\begin{align*}
\vect{s}_{\text{roof}} &= (0.25/0.02,\ -0.08/0.04,\ 0.10/0.01)=(12.5,\ -2,\ 10), & \vect{s}\cdot\vect{s}&=260.25\\
\vect{s}_{\text{soil}} &= (0.15/0.02,\ -0.15/0.04,\ 0.05/0.01)=(7.5,\ -3.75,\ 5), & \vect{s}\cdot\vect{s}&=95.3125 .
\end{align*}
A 20\% roof: $\vect{d}=0.2\,\vect{s}_{\text{roof}}=(2.5,\ -0.4,\ 2)$, $\vect{d}\cdot\vect{d}=6.25+0.16+4=10.41$.
\begin{align*}
\text{soil: } \vect{s}\cdot\vect{d} &= 18.75+1.5+10=30.25,\\
\text{ACE} &= 30.25^2/(95.3125\times10.41)=915.06/992.20=0.9223 .
\end{align*}
Roof scores 1.00 (the directions are parallel). A 50\% roof: $\vect{d}=(6.25,\ -1,\ 5)$, $\vect{d}\cdot\vect{d}=65.0625$,
soil $\vect{s}\cdot\vect{d}=75.625$, ACE $=5719.14/6201.27=0.9223$, the same. Scaling $\vect{d}$ by the fill $a$
multiplies the numerator by $a^2$ and the denominator by $a^2$: \textbf{fill changes the length of $\vect{d}$, not its
direction}.

Noise does not scale with fill. Add one standard deviation per band, whitened $(1,1,-1)$. At fill 0.2,
$\vect{d}=(3.5,\ 0.6,\ 1)$, roof $\vect{s}\cdot\vect{d}=52.55$, $\vect{d}\cdot\vect{d}=13.61$, ACE $=2761.50/3542.00=0.780$.
At fill 0.5, $\vect{d}=(7.25,\ 0,\ 4)$, $\vect{s}\cdot\vect{d}=130.625$, $\vect{d}\cdot\vect{d}=68.5625$, ACE
$=17062.89/17843.39=0.956$ \tagC.
\end{worked}

\paragraph{Leave one spectrum out.} To measure naming, the check plants each library spectrum and matches it against
the other 410, setting its own score to $-1$ (\code{ace[e] = -1}); matching a spectrum against itself would always
win and measure nothing.

\subsection{Matching versus a classifier}
\index{dual encoder}

A \emph{closed classifier} ends in one output per class; a new class needs a new output and retraining. A
\emph{matching model} (a dual encoder) has two networks: one embeds the pixel and its ring as a short vector, the
other embeds a library spectrum as a vector of the same length; the score is their cosine, and naming takes the
best library entry. A new library needs no retraining: embed its spectra and they are candidates. With a pixel
embedded at $(0.8,0.6)$ and library entries painted metal $(1,0)$, tarp $(0,1)$ and a newly added glass spectrum
$(0.6,0.8)$, the cosines are $0.80$, $0.60$ and $0.8\times0.6+0.6\times0.8=0.96$: the new entry wins without any
training \tagA. ACE is itself a matching model whose two ``encoders'' are fixed: subtract the ring and whiten.
\textbf{Level~1 does not depend on any library}, because it models nature; a new library changes only level~2's
candidates. The price is that a matching model generalises only to spectra like those its encoder learned from,
while a classifier trained on a fixed list can be sharper on that list. The companion volume's Part~VII derives the
contrastive loss that trains a dual encoder; \cref{ch:namer} shows the class-level version that was built.

\subsection{The separability check}\label{sec:material:separability}

Before training anything, \file{separability.py} bounds what a strong classical matcher can do: ``No training - this
bounds what the physics and the library allow before we design classes.'' For each scene (the four PRISMA bench
scenes; three implant-bench host scenes each for EnMAP and AVIRIS-NG) it resamples the 411 spectra to the scene's
bands, plants each into 12 real background pixels at least 8 pixels from any object as
$\vect{x}=g\,(a\vect{t}+(1-a)\vect{b})+\vect{n}$, with a smooth gain error $|g-1|\le0.1$ and the scene's own per-band
noise, and matches $\vect{d}=\vect{x}-g\,\vect{m}$ by ACE against every other spectrum, whitened with the scene's
covariance (30 components plus a ridge). The best match's family is level~2, its group level~1. Plantings per fill:
PRISMA $411\times12\times4=19{,}728$; EnMAP and AVIRIS-NG $411\times12\times3=14{,}796$.

\begin{table}[htbp]
\centering\small
\begin{tabular}{ccccccc}
\toprule
 & \multicolumn{3}{c}{L1: man-made right / natural right} & \multicolumn{3}{c}{L2: class right}\\
\cmidrule(lr){2-4}\cmidrule(lr){5-7}
fill & PRISMA & EnMAP & AVIRIS-NG & PRISMA & EnMAP & AVIRIS-NG\\
\midrule
0.10 & 0.61 / 0.78 & 0.48 / 0.72 & 0.59 / 0.77 & 0.38 & 0.23 & 0.36\\
0.25 & 0.77 / 0.83 & 0.66 / 0.79 & 0.76 / 0.81 & 0.55 & 0.44 & 0.53\\
0.50 & 0.85 / 0.86 & 0.80 / 0.84 & 0.84 / 0.85 & 0.61 & 0.58 & 0.59\\
1.00 & 0.87 / 0.86 & 0.87 / 0.87 & 0.87 / 0.85 & 0.64 & 0.64 & 0.62\\
\bottomrule
\end{tabular}
\caption{The separability check, all three sensors, final on 2026-10-31. On PRISMA, level~2 right \emph{given} level~1
right is 0.63 / 0.72 / 0.72 / 0.73. Source: \tagM{research/experiments/spectral-namer/README.md} and the three
\texttt{separability\_<sensor>.json} files.}
\label{tab:material:separability}
\end{table}

A class can be right only if level~1 was, so \cref{eq:material:twolevel} holds on the counts: at PRISMA fill 0.10,
$0.6096\times0.6309=0.3846$, the table's 0.38 \tagC.

Per family, PRISMA at fill 0.25 (from the JSON's confusion counts \tagC): oil $1537/1920=80\%$, plastic
$1319/1872=70\%$, textile $464/672=69\%$, wood and paper $239/384=62\%$, paint $174/384=45\%$, metal $255/624=41\%$,
building $205/528=39\%$ (16\% called pavement), roofing $398/1296=31\%$, pavement $75/720=10\%$ (15\% called soil,
13\% rock, 12\% building), glass $0/48$.

\begin{sourcenote}[: pavement is 10\%, not 12\%]
The experiment README gives pavement at fill 0.25 as ``$\sim$12\%''. The JSON's pavement row has 75 of 720 right,
10.4\%; the 12\% is pavement called \emph{building}, 87 of 720.
\end{sourcenote}

Four readings. Level~1 is the harder half at small fills: at 10\% fill, 39\% of man-made plantings are called
natural on PRISMA and 52\% on EnMAP, and most targets at 30\,m are that small; at full fill all three sensors reach
0.87. Materials with their own absorption features (hydrocarbon, polymer C-H, cellulose) separate well: oil,
plastic, textile, wood and paper. The use-based families fail as the regex section predicted: roofing is scattered,
building and pavement confuse each other and natural rock and soil, and glass, with one spectrum, can never be
matched once it is left out. And these are library spectra planted in real pixels, so the laboratory-to-satellite
gap is absent: a real namer will do worse.

The README turns this into level~2's class list: \emph{oil/fuel, plastic, textile/fabric, wood/paper,
paint/coatings, metal, concrete/asphalt} (building and pavement merged, roofing dissolved into its materials), with
glass/PV once spectra exist and camouflage and military coatings once field spectra arrive. ``These classical
numbers are the baseline the learned namer must beat, mainly at 10-25\% fill.'' This list of \emph{materials} is
narrower than the analyst's list of \emph{things} above: vessels, disturbed earth and cut vegetation are contexts
or mixtures, not single spectra.

\begin{keypoint}
Name what nature cannot explain before naming what it is, multiply the two answers, and match against spectra
rather than classify into a fixed list, so that the unknown is flagged and a new library costs no retraining.
\end{keypoint}

\begin{drift}
The design (2026-10-31, now \file{docs/proposals/spectral-namer.md}), the class lists, the family counts (411 spectra,
cap 40), and every separability number; a new library, class list or set of host scenes changes them.
\end{drift}

%======================================================================
\section{Export: where geography comes back}\label{sec:material:export}

\subsection{Every action GeoTIFF has an identity transform}

A GeoTIFF places its pixel grid on the Earth through an \emph{affine transform}\index{affine transform}: six
numbers $(a,b,c,d,e,f)$ that map a column $j$ and row $i$ to map coordinates,
\begin{equation}
X=a\,j+b\,i+c,\qquad Y=d\,j+e\,i+f,
\label{eq:material:affine}
\end{equation}
together with a \emph{coordinate reference system}\index{coordinate reference system} (CRS) that says what
$X$ and $Y$ mean (degrees of longitude and latitude, metres in a UTM zone, \dots). Vendables carry no spatial
reference (\cref{ch:product}), so every raster an action writes uses the \emph{identity} transform,
$a=e=1$ and $b=c=d=f=0$, and no CRS (\file{_anomaly_scoring_render.py}, line 61)\index{identity transform}:
\[
X=j,\qquad Y=i.
\]
Pixel coordinates pretending to be map coordinates. Analysis works in pixel grids and does not need more;
only the deliverable does.

\subsection{Recovering the CRS at export}

The export endpoint, \code{POST /actions/\{id\}/export} in \file{backend/allotrope/api/action_export.py},
runs synchronously in the api. It re-opens the raw scene file through \file{app/georef/from_scene_file.py},
which knows each sensor:
\begin{itemize}
\item \textbf{EnMAP L2A, AVIRIS-NG, Landsat~9, HotSat}: the raw file is already a georeferenced raster;
  rasterio reads its CRS and transform. If the vendable's shape differs from the source's, the transform is
  scaled so that the outer corners stay pinned.
\item \textbf{PRISMA L2D}: since 2026-10-31 the product's own UTM grid, from the HE5's root attributes
  \code{Epsg_Code} and \code{Product_{UL,UR,LL}corner_{easting,northing}}, which are the corner pixels'
  \emph{centres} (\code{prisma_utm_grid}; worked for Set-1 in \cref{we:submission:set1}). Only if those are
  missing, or the vendable's shape differs from the product's, does it fall back to the older method: the minimum and
  maximum of the per-pixel latitude and longitude arrays, as an axis-aligned EPSG:4326 transform built with
  rasterio's \code{from_bounds}.
\end{itemize}
If none can be recovered the export returns \code{422 crs_missing} rather than a bundle that a downstream
evaluation would disqualify for having no projection.

\begin{sourcenote}[: the bounding box was not a fit, and it was pixels off]
The anomaly course used to say export ``reads the per-pixel latitude and longitude arrays \dots\ and fits an affine
transform'' (AD13). The code never fitted: it took $\min$ and $\max$ of the arrays and called
\code{from_bounds(min_lon, min_lat, max_lon, max_lat, W, H)}, which is exact only for a north-aligned swath and treats
pixel-centre coordinates as edges. Its docstring called that ``the industry-standard simplification for L2D HCO data''.
Measured on Set-1, it put pixels a median 2.4\,px (95\% 4.1, max 5.6) from the product grid
\tagM{docs/proposals/aigc-evaluation-submission.md}; the evaluation, which compares masks in pixel space, was never
affected. Since 2026-10-31 it is only the fallback, and the docstring now says ``approximate, can be pixels off''.
\end{sourcenote}

The worked example below is the fallback, kept because it shows what an affine transform does; the product grid is
worked in \cref{we:submission:set1}.

\begin{worked}[label={we:material:georef}]{An axis-aligned transform from a bounding box (the fallback)}
A PRISMA frame of $W=1219$ columns and $H=1210$ rows whose latitude and longitude arrays span
$77.00^\circ$ to $77.36^\circ$ east and $12.80^\circ$ to $13.13^\circ$ north (an assumed box, for
illustration \tagA). \code{from_bounds} sets
\begin{align*}
a &= \frac{\text{east}-\text{west}}{W} = \frac{77.36-77.00}{1219} = \frac{0.36}{1219} = 0.00029532^\circ,\\
e &= -\frac{\text{north}-\text{south}}{H} = -\frac{13.13-12.80}{1210} = -\frac{0.33}{1210} = -0.00027273^\circ,\\
c &= \text{west} = 77.00, \qquad f = \text{north} = 13.13, \qquad b=d=0.
\end{align*}
$e$ is negative because rows count downwards while latitude grows northwards. The centre of the pixel in row
605, column 610 is at $j=610.5$, $i=605.5$:
\begin{align*}
X &= 0.00029532\times610.5+77.00 = 0.18029+77.00 = 77.18029^\circ\text{ E},\\
Y &= -0.00027273\times605.5+13.13 = -0.16514+13.13 = 12.96486^\circ\text{ N}.
\end{align*}
How big is a pixel on the ground? A degree of latitude is about 110.6~km; a degree of longitude at
$13^\circ$ north is about $111.32\times\cos13^\circ=111.32\times0.9744=108.47$~km. So
\begin{align*}
\text{width} &\approx 0.00029532\times108{,}470\text{ m} = 32.0\text{ m},\\
\text{height} &\approx 0.00027273\times110{,}574\text{ m} = 30.2\text{ m},
\end{align*}
close to PRISMA's 30~m, as it should be. The book's script reproduces these with rasterio \tagC.
\end{worked}

\subsection{What the bundle holds}
\index{export bundle}

For hyperspectral scenes the zip is built by \file{app/spectral_match/export.py}. Every name starts with
\texttt{hyper\_}, as the submission rules require, and the manifest records the schema:
\texttt{allotrope-export-hyper/2} for actions with family matching, \texttt{/1} for older SAM-only ones.
\begin{lstlisting}[language={}]
hyper_<scene8>_<action8>/
    hyper_anomaly_materials.tif      int16 family codes, nodata -2, real CRS and transform
    hyper_anomaly_materials.shp      one point per anomaly pixel (+ .shx .dbf .prj .cpg)
                                     DBF: family, group, fconf, named, best_ent, mat1..3, scr1..3
    hyper_material_legend.json       code -> family, group, colour; alpha and p*
    hyper_sam_legend.json            SAM's library index -> material
    hyper_match_summary.csv          flat table, no geometry
    MANIFEST.json                    schema allotrope-export-hyper/2
\end{lstlisting}
Material identity travels three ways (a GeoTIFF metadata tag, the DBF columns, the legend file) because
different GIS tools read different parts, and a shortlist nobody can open is worthless. A committed
\code{anomaly_detection_prep} has its own bundle, from \file{app/anomaly_detection/export.py}: a binary 0/1 mask and
one point per anomaly pixel. Until 2026-10-31 it was always named \texttt{thermal\_}; it is now \texttt{thermal\_}
for Landsat~9 and HotSat-1 and \texttt{hyper\_} otherwise, with no score attribute in its shapefile. That bundle, not
the materials bundle above, is what an evaluation submission uses (\cref{ch:submission}).

%======================================================================
\section{Using Indradhanu V2 in an action chain}\label{sec:material:v2}

\subsection{Two words, and what V2 is}

A \emph{codename}\index{codename} is the name you pick a model by in the scoring action: ``Indradhanu'',
``MNF-RX'', ``Indradhanu V2''. The resolver maps it to a checkpoint and its capabilities
(\file{backend/allotrope/foundation_models/resolver.py}). \emph{Native bands}\index{native bands} are the
wavelengths the sensor actually measured, PRISMA's own list or EnMAP's own list, as opposed to the 165-band
common grid that \code{band_filter_apply} resampled to by default until 2026-10-31 (\cref{ch:bandpipeline}).

Indradhanu~V2 is the spectral-stem model of Part~IV, fine-tuned with a segmentation head (\cref{ch:finetune})
and shipped as an opt-in codename (\cref{ch:shipping}). It reads native bands, scores with residual
Mahalanobis like Indradhanu, and sets to 0 the pixels its own head calls cloud.

\begin{drift}
Indradhanu~V2 was wired into \code{anomaly_scoring} on 2026-10-19 serving FT04 (\texttt{v2.0.0}); on
2026-10-29 it was switched to Indradhanu Heavy, FTH3 epoch~5 (\texttt{v2.1.0}, 35.3~M parameters)
(\file{final design/ROADMAP.md}, row OB-14; \file{docs/04-models.md}). On the real PRISMA bench with cloud
exclusion: v0.2.0 pAUC 0.165, FT04 0.170 (a tie), Heavy 0.189 ($+0.020$ over FT04, resolved). CPU cost on the
worker: 1819~s per PRISMA scene for Heavy against 1058~s for FT04 on 2 threads (ROADMAP OB-14); 454~s and 274~s
since the compose file raised the worker to 16 threads on 2026-10-29 (\cref{we:heavy:threads}). The scoring codename
defaults still pick the old ``Indradhanu'' (v0.2.0); since 2026-10-31 the band-filter defaults keep native bands
(\cref{ch:submission}).
\end{drift}

\begin{analogy}
A film lab with two developers. The old one only takes 35~mm film; the new one only takes medium format. You
choose the format when you load the camera, long before you reach the lab. Hand the new developer a 35~mm
roll and it is refused at the door with a note: ``medium format only, reload the camera''. And one trip to
the lab cannot develop both, because one camera was loaded with one kind of film. The camera is
\code{band_filter_apply}; the developers are the two codenames; the note is the error below.
\end{analogy}

\subsection{The chain, setting by setting}

Two settings matter; since 2026-10-31 the first is the sensor default, so in the UI only the second is a choice:
\begin{lstlisting}[language={}]
Scene -> Project
  band_filter_apply        use_common_wavelength_grid: false     <- change 1 (default since 2026-10-31)
  anomaly_scoring          model_codenames: ["Indradhanu V2"]    <- change 2
  anomaly_detection_prep   unchanged: reads models/indradhanu_v2/anomaly_score.tif
  spectral_library_match   unchanged: reads the native onboarding cube anyway
\end{lstlisting}
The per-model folder is the codename in lower case with spaces turned into underscores (\code{_safe_dirname}):
\begin{align*}
\text{``Indradhanu V2''}\ \to\ \text{``indradhanu v2''}\ \to\ \text{``indradhanu\_v2''}.
\end{align*}

\subsection{The error you get otherwise}
\index{band-grid check}

Give V2 a band filter on the 165-band grid (before 2026-10-31 the default; now \code{use_common_wavelength_grid: true},
or an API call that omits the field, whose pydantic default is still \code{true}) and the action is still accepted
(the api cannot see the upstream band filter's settings when you submit). The worker starts and, before loading the
model, runs \code{_check_band_grid} (\file{_anomaly_scoring_run.py}). The action fails with:
\begin{lstlisting}[language={}]
model 'Indradhanu V2' reads native bands: run band_filter_apply with
use_common_wavelength_grid=false (this one resampled to the 165-band grid)
\end{lstlisting}
The reverse, the old ``Indradhanu'' on a native band filter, used to fail with ``needs the 165-band grid: run
\code{band_filter_apply} with \code{use_common_wavelength_grid=true}''. Since 2026-10-31 it runs: the worker resamples
the cube to the grid for that model alone (\code{_to_common_grid}, ``the same resampler \code{band_filter_apply} uses''), and
a comment where the check was says common165 on native bands ``is not an error''. Before 2026-10-19 nothing checked
any of this.

How does the worker know which grid it was given? Not by counting bands:
\begin{lstlisting}
on_grid = len(wavelengths) == len(grid) and np.allclose(wavelengths, grid)   # grid: 460, 470, ..., 2450 nm
\end{lstlisting}
A native list with exactly 165 bands would still not equal $460, 470, \dots, 2450$~nm and would count as
native. Classical detectors skip the check and run on either grid; so do thermal scenes.

\subsection{Running both models}

\code{anomaly_scoring} takes one upstream band filter (\code{input_band_filter_output_id}), and every model in
\code{model_codenames} reads that one output. Since 2026-10-31, on a native band filter, one action can run both: V2
reads the native cube and v0.2.0 gets it resampled. Whether v0.2.0 on a resampled native cube matches v0.2.0 on a
165-grid band filter exactly has not been measured; the resampler is the same function. Before 2026-10-31, V2 needed
native bands and v0.2.0 the grid, whichever did not match failed the whole action, and comparing them took two
branches in the same Project:

\begin{figure}[htbp]\centering
\begin{tikzpicture}[node distance=5mm and 7mm,
  box/.style={draw,rounded corners,fill=#1!8,draw=#1!70!black,minimum height=8mm,align=center,font=\small}]
\node[box=maskgrey] (s) {Scene $\to$ Project};
\node[box=spec,above right=2mm and 8mm of s] (b1) {\texttt{band\_filter\_apply}\\grid: true};
\node[box=spec,below right=2mm and 8mm of s] (b2) {\texttt{band\_filter\_apply}\\grid: false};
\node[box=attn,right=of b1] (a1) {\texttt{anomaly\_scoring}\\{}[``Indradhanu'']};
\node[box=attn,right=of b2] (a2) {\texttt{anomaly\_scoring}\\{}[``Indradhanu V2'']};
\draw[-Stealth] (s.east) -- (b1.west);
\draw[-Stealth] (s.east) -- (b2.west);
\draw[-Stealth] (b1) -- (a1);
\draw[-Stealth] (b2) -- (a2);
\end{tikzpicture}
\caption{Before 2026-10-31, comparing the two Indradhanu models on one scene took two branches, each with its
own band filter and scoring action. Since then one native band filter and one scoring action serve both. A
classical detector can ride along in either.}\label{fig:material:branches}
\end{figure}

One consequence of the old limit: \code{anomaly_detection_prep} reads a single scoring output
(\code{input_anomaly_scoring_output_id}), so V2 and v0.2.0 could not be weighted into one composite. With both in one
action on a native band filter, they can.

\subsection{What comes out, and how to read it}

Under the scoring action's output, in \file{models/indradhanu_v2/}, next to the usual files
(\file{_anomaly_scoring_run.py}, lines 471--491 and 595--601):
\begin{center}\small
\begin{tabular}{@{}lp{0.62\linewidth}@{}}
\toprule
output & contents\\
\midrule
\texttt{mask\_cloud.tif} & uint8, 1 where the score was set to 0 (V2 called the pixel cloud)\\
\texttt{class\_map.tif} & uint8, V2's classes: 0 land, 1 water, 2 cloud, 3 shadow, 4 snow\\
\texttt{stats.json}: \texttt{bands\_read} & how many bands V2 read (bands it never trained on are dropped first)\\
\texttt{stats.json}: \texttt{excluded\_classes} & \texttt{["cloud"]}\\
\texttt{stats.json}: \texttt{cloud\_rule} & \texttt{"thick"} or \texttt{"cloud\_or\_thin"}\\
\texttt{stats.json}: \texttt{excluded\_percent} & $100\times$ cloud pixels / spatially valid pixels, to 3 places\\
\bottomrule
\end{tabular}
\end{center}
The class map is built from the segmentation head's five sigmoid outputs, averaged over $128\times128$ windows
at stride 64, by priority: cloud over snow over shadow over water over land
(\file{app/foundation_models/inferencers/spectral_stem_inferencer.py}, lines 112--141). Only cloud is
excluded; water is in the class map but \emph{not} zeroed, because vessels on water are in scope.

\begin{sourcenote}[: since 2026-10-29 ``cloud'' means thick cloud only]
The anomaly course (AD23, written for FT04) says \texttt{mask\_cloud.tif} marks pixels V2 ``called cloud or
thin cloud''. Since V2 became Indradhanu Heavy, the resolver sets \code{thin_cloud_is_cloud=False}
(\file{backend/allotrope/foundation_models/resolver.py}, line 176): Heavy's thin-cloud output fires on haze, so
only thick cloud is zeroed, and the new stats key \texttt{cloud\_rule} records \texttt{"thick"}
(\file{_anomaly_scoring_run.py}, lines 471--481).
\end{sourcenote}

\paragraph{The order matters.}\index{cloud exclusion} V2 scores every pixel in the keep mask first and sets cloud to 0
\emph{afterwards}. A cloud pixel still takes part in fitting the residual Mahalanobis statistics of
\cref{ch:residualmaha}; only its own score becomes 0. That is the semantics the evaluator measured. A pixel
scored 0 is not ``unknown'': it sinks to the bottom of the ranking. Read \texttt{excluded\_percent} as the share
of the scene in which you cannot find anything.

\begin{worked}[label={we:material:excluded}]{The excluded percentage}
A scene of $1000\times1200$ pixels, $72{,}000$ invalid, $90{,}240$ called cloud (AD23's example \tagA):
\begin{align*}
\text{total} &= 1000\times1200 = 1{,}200{,}000,\\
\text{valid} &= 1{,}200{,}000-72{,}000 = 1{,}128{,}000,\\
\texttt{excluded\_percent} &= \frac{100\times90{,}240}{1{,}128{,}000} = \frac{9{,}024{,}000}{1{,}128{,}000} = 8.000.
\end{align*}
\end{worked}

\paragraph{Two caveats.} EnMAP is cleaned twice: its default band filter applies the provider's quality masks
(cloud, cloud shadow, haze) before scoring, and V2 then zeros its own cloud on top; this is the product's
convention, but it is not how V2's PRISMA numbers were measured. And native PRISMA bands run downwards in
wavelength (\cref{ch:sortbug}); \code{scene_segmentation} now picks its visible bands by wavelength, not by
position, so a native-band chain gets the right brightness too. The viewer does not show V2's masks yet; they
are on disk only.

%======================================================================
\begin{confusions}
\confusion{Why not a fixed threshold on residual Mahalanobis, which has a chi-squared-like scale?}{Its scale
depends on the scene's covariance, the number of kept directions and the local ring; the tail is heavy and
scene-specific. A percentile adapts; a fixed number does not.}
\confusion{Top 0.5 percent means 0.5 percent of the kept pixels.}{Only if a classical detector's NaNs remove
the unkept pixels from the composite. With foundation models alone, the unkept pixels are finite zeros and the
percentile is taken over the whole frame.}
\confusion{A confident family name means the pixel is that material.}{No. The confidence is a percentile
against the scene's ordinary pixels, and false alarms are unusual by definition. It is a lead for a person.}
\confusion{ACE ignores the fill, so the fill does not matter for naming.}{ACE's formula ignores it; its accuracy
does not, because noise stays the same size while the signal shrinks with the fill
(\cref{we:material:ace}, \cref{tab:material:separability}).}
\confusion{The two-level namer and the companion volume's two-level covariance are the same idea.}{No. The
two-level covariance is a model of noise (a few loud directions, many quiet ones). The two-level namer is two
questions asked in turn.}
\confusion{The export can be georeferenced because the score map is.}{The score map is not: every action
raster has an identity transform. The CRS is recovered from the raw file at export, or the export fails with
422.}
\confusion{The error from V2 means the model is broken.}{It means the band filter is on the wrong grid. The
message names the setting to change.}
\confusion{One scoring action cannot run both Indradhanus.}{It could not before 2026-10-31. On a native band filter it
now can: v0.2.0's cube is resampled inside scoring.}
\confusion{\texttt{mask\_cloud.tif} is the keep mask.}{No. The keep mask comes from
\texttt{scene\_segmentation}'s thresholds and decides which pixels the statistics are fitted on.
\texttt{mask\_cloud.tif} comes from V2's own head and is applied after scoring.}
\confusion{\texttt{spectral\_library\_match} needs the native band filter now.}{It never used the band filter's
output; it reads the native onboarding cube.}
\end{confusions}

\begin{checkyourself}
\item\label{ex:material:1} A scene has $2{,}400{,}000$ finite composite pixels with all values different. How
  many does $q=99.75$ flag? Use \cref{eq:material:nflag} and show every step.
\item\label{ex:material:2} Two models with weights 1 and 3. At a pixel the first reads 0.4 on a range
  $[0.1,1.6]$, the second 12 on $[0,60]$. Compute the composite value.
\item\label{ex:material:3} With $\alpha=0.02$ and 15 man-made families, compute $p^*$, and the bound on the
  chance that an ordinary pixel gets some man-made name.
\item\label{ex:material:4} A GeoTIFF has $a=30$, $b=0$, $c=500{,}000$, $d=0$, $e=-30$, $f=1{,}400{,}000$ in a
  UTM zone (metres). Where is the centre of the pixel in row 10, column 20? What would the identity transform
  give for the same pixel?
\item\label{ex:material:5} A PRISMA Project has a band filter created with \code{use_common_wavelength_grid: true}
  and an \code{anomaly_scoring} action with \code{model_codenames: ["Indradhanu V2", "MNF-RX"]}. What happens, and
  which single setting fixes it? What happens with today's default band filter?
\item\label{ex:material:6} A scene is $800\times900$ pixels, $20{,}000$ invalid, and V2 calls $35{,}800$ valid
  pixels cloud. Compute \texttt{excluded\_percent}. Is a cloud-shadowed vessel's score zero?
\item\label{ex:material:7} Level~1 gives $P(M)=0.7$ and level~2 gives bare metal 0.5, painted metal 0.3, tarp 0.2.
  Compute every class's probability and $P(\text{natural})$, and check that they add to 1.
\item\label{ex:material:8} In \cref{we:material:ace}, make the pixel 100\% roof with no noise. Compute the whitened
  $\vect{d}$ and ACE against soil, and say why the answer could have been written down without arithmetic. Then, from
  \cref{tab:material:separability} at fill 0.25, compute $P(\text{L1 right})\times P(\text{L2 right}\mid\text{L1
  right})$ from the unrounded 0.7724 and 0.7151.
\end{checkyourself}

\nextchapter{ch:measures}
